% from ega4-11-fr.tex
\oldpage[IV-3]{151}
% IV3_B22_P151_ANCHOR
l'homomorphisme canonique
(\textbf I, 2.4.4)~; on suppose que les intersections finies des idéaux
$u_\alpha^{-1}(\mathfrak m_\alpha'^{\,n_\alpha})$ forment un système fondamental
de voisinages de $0$ dans $\mathcal O_y$. Posons
$X'_\alpha=X\times_Y Y'_\alpha$, $f'_\alpha=f_{(Y'_\alpha)}$,
$\mathcal F'_\alpha=\mathcal F\otimes_Y Y'_\alpha$, et supposons vérifiée l'une
des hypothèses suivantes~:
\begin{enumerate}[label=(\roman*)]
  \item Pour tout $\alpha$, $\mathcal F'_\alpha$ est $f'_\alpha$-plat en tous
  les points dont la projection dans $X$ est égale à $x$ et la projection dans
  $Y'_\alpha$ égale à $y'_\alpha$.
  \item Pour tout $\alpha$, il existe un $x'_\alpha\in X'_\alpha$ dont la
  projection sur $X$ est $x$ et la projection dans $Y'_\alpha$ égale à
  $y'_\alpha$, tel que $\mathcal F'_\alpha$ soit $f'_\alpha$-plat au point
  $x'_\alpha$, et $k(x)$ est une extension primaire de $k(y)$.
\end{enumerate}
Dans ces conditions, $\mathcal F$ est $f$-plat au point $x$.
\end{theorem}

\begin{proof}
Soit $\mathfrak m_y$ l'idéal maximal de $\mathcal O_y$~; comme $\mathcal O_y$
et $\mathcal O_x$ sont noethériens et $\mathcal O_y\to\mathcal O_x$ un
homomorphisme local, il suffit, en vertu de
$(\mathbf 0_{\mathrm{III}},10.2.2)$, de prouver que pour tout $n>0$,
$\mathcal F_x/\mathfrak m_y^n\mathcal F_x$ est un
$(\mathcal O_y/\mathfrak m_y^n)$-module plat. Désignons par
$(\mathfrak J_\lambda)$ la famille des intersections finies des
$u_\alpha^{-1}(\mathfrak m_\alpha'^{\,n_\alpha})$~; par hypothèse, il existe
$\lambda$ tel que $\mathfrak J_\lambda\subset\mathfrak m_y^n$, et comme
\[
\mathcal F_x\otimes_{\mathcal O_y}(\mathcal O_y/\mathfrak m_y^n)
=\bigl(\mathcal F_x\otimes_{\mathcal O_y}(\mathcal O_y/\mathfrak J_\lambda)\bigr)
 \otimes_{\mathcal O_y/\mathfrak J_\lambda}(\mathcal O_y/\mathfrak m_y^n),
\]
il suffira de prouver que $\mathcal F_x/\mathfrak J_\lambda\mathcal F_x$ est un
$(\mathcal O_y/\mathfrak J_\lambda)$-module plat. Or, pour chaque $\alpha$ tel
que $\mathfrak J_\lambda\subset u_\alpha^{-1}(\mathfrak m_\alpha'^{\,n_\alpha})$,
on a, par passage aux quotients, un homomorphisme
$u'_\alpha:\mathcal O_y/\mathfrak J_\lambda\to
A'_\alpha/\mathfrak m_\alpha'^{\,n_\alpha}$, et l'intersection des noyaux des
$u'_\alpha$ est réduite à $0$. Compte tenu de (\textbf I, 3.6.1), on voit qu'on
est ramené à (11.4.11) dans le cas (ii), et à (11.4.12) dans le cas (i).
\end{proof}

\begin{corollary}[11.5.2]
\label{IV.11.5.2-fr}
Soient $Y$ un préschéma localement noethérien, $f:X\to Y$ un morphisme
localement de type fini, $\mathcal F$ un $\mathcal O_X$-Module cohérent, $x$ un
point de $X$, $y=f(x)$~; posons $A=\mathcal O_y$. Soit $u:A\to A'$ un
homomorphisme de $A$ dans un anneau de Zariski $A'$, tel que l'image réciproque
$u^{-1}(\mathfrak r')$ du radical $\mathfrak r'$ de $A'$ soit l'idéal maximal
$\mathfrak m$ de $A$~; supposons en outre que l'homomorphisme
$\widehat u:\widehat A\to\widehat{A'}$ soit injectif. Posons
$Y'=\operatorname{Spec}(A')$, $X'=X\times_Y Y'$, $f'=f_{(Y')}$,
$\mathcal F'=\mathcal F\otimes_Y Y'$. Pour que $\mathcal F$ soit $f$-plat au
point $x$, il faut et il suffit que $\mathcal F'$ soit $f'$-plat en tout point
dont la projection dans $X$ est égale à $x$ et dont la projection dans $Y'$ est
égale à un point fermé de $Y'$.

Si en outre $A'$ est une $A$-algèbre finie, on peut dans ce qui précède
remplacer l'hypothèse que $\widehat u$ est injectif par l'hypothèse que $u$ est
injectif.
\end{corollary}

\begin{proof}

Comme $A$ (resp. $A'$) s'identifie à un sous-anneau de $\widehat A$ (resp.
$\widehat{A'}$) (Bourbaki, \emph{Alg. comm.}, chap.~III, \S~3, n$^\circ$~3,
prop.~6), on voit d'abord que $u$ lui-même est injectif et que $\widehat u$ est
son prolongement par continuité à $\widehat A$.

Soit $(\mathfrak m'_\alpha)$ la famille des idéaux maximaux de $A'$~; comme on a
\[
\mathfrak m_\alpha'^{\,n}\widehat{A'}
=\widehat{\mathfrak m}_\alpha'^{\,n},
\qquad\text{et}\qquad
\widehat{\mathfrak m}_\alpha'^{\,n}\cap A'
=\mathfrak m_\alpha'^{\,n},
\]
et que les $\mathfrak m_\alpha'^{\,n}$ sont ouverts dans $A'$, on a
$u^{-1}(\mathfrak m_\alpha'^{\,n})
=A\cap\widehat u^{-1}(\widehat{\mathfrak m}_\alpha'^{\,n})$, et il suffira de
montrer que dans $\widehat A$, les intersections finies des
$\widehat u^{-1}(\widehat{\mathfrak m}_\alpha'^{\,n_\alpha})$ forment un système
fondamental de voisinages de $0$, ce qui permettra d'appliquer (11.5.1) aux
homomorphismes composés $v_\alpha\circ u:A\to A'_{\mathfrak m'_\alpha}$,
où $v_\alpha:A'\to A'_{\mathfrak m'_\alpha}$ est l'homomorphisme canonique.
Comme $\widehat A$ est complet, il suffira de montrer que l'intersection des
$\widehat u^{-1}(\widehat{\mathfrak m}_\alpha'^{\,n_\alpha})$ est réduite à
$0$ (Bourbaki, \emph{Alg. comm.}, chap.~III, \S~2, n$^\circ$~7, prop.~8, où on
peut dans la démonstration remplacer la suite décroissante par un ensemble
filtrant quelconque). Or, pour tout $\alpha$
